% ============================================================
% f2014_liaoning_science_q19_tetrahedron.tex
% 2014 辽宁卷（理科）第 19 题：两个平面垂直的四面体
% 来源 PDF：普通高考/2014/2014辽宁理.pdf
% 原页：第 2 页（A3），印面 896，第 19 题题干下方
% 图中 E、F 分别为 AC、DC 的中点；AB、BC、BD、BE 为被遮住的棱，
% 画虚线。
% ============================================================
\begin{tikzpicture}[>=Stealth, line width=0.7pt, font=\small, x=1cm, y=1cm,
  every node/.style={inner sep=1.5pt}]
  \coordinate (A) at (2.00,4.00);
  \coordinate (C) at (4.90,2.20);
  \coordinate (D) at (1.35,0.35);
  \coordinate (B) at (2.75,1.95);
  \coordinate (E) at (3.45,3.10);
  \coordinate (F) at (3.125,1.275);

  % 可见棱与截面线
  \draw (A) -- (C) -- (D) -- cycle;
  \draw (E) -- (F);

  % 被遮住的棱
  \draw[dashed] (A) -- (B);
  \draw[dashed] (B) -- (C);
  \draw[dashed] (B) -- (D);
  \draw[dashed] (B) -- (E);

  % 顶点与中点
  \fill (A) circle (1.1pt) node[above] {$A$};
  \fill (C) circle (1.1pt) node[right] {$C$};
  \fill (D) circle (1.1pt) node[below left] {$D$};
  \fill (B) circle (1.1pt) node[below right] {$B$};
  \fill (E) circle (1.1pt) node[above right] {$E$};
  \fill (F) circle (1.1pt) node[below right] {$F$};
\end{tikzpicture}
